% Source: https://artofproblemsolving.com/community/c3h1847387_solution
% Stable identity: aops-topic-1847387-q01
Since $5$ and $-6$ are the roots, $y=a(x-5)(x+6)$ for some $a\ne0$. At $x=0$: $a(-5)(6)=-30a=-1$, so $a=\dfrac{1}{30}$. Hence
\[
y=\dfrac{(x-5)(x+6)}{30}=\boxed{\dfrac{1}{30}x^2+\dfrac{1}{30}x-1}.
\]

\begin{asy}
usepackage("amsmath");
size(6cm);
defaultpen(fontsize(10pt));
import graph;
size(8cm,5cm,IgnoreAspect);
real f(real x){return (x^2+x-30)/30;}
draw((-8.6,0)--(8,0),gray(0.6),Arrow(5));
draw((0,-1.35)--(0,1.15),gray(0.6),Arrow(5));
draw(graph(f,-8.2,7.2),blue+linewidth(0.9));
dot((-6,0));dot((5,0));dot((0,-1));
label("$(-6,0)$",(-6,0),SW);label("$(5,0)$",(5,0),SE);label("$(0,-1)$",(0,-1),SE);
label("$x$",(8,0),E);label("$y$",(0,1.15),N);
label("\mbox{}",(min(currentpicture,user=true)+max(currentpicture,user=true))/2,black); // keeps the text after the figure black when LaTeX typesets the labels inline
\end{asy}

(The graph is drawn with a stretched vertical scale.)
